\section{Results}

\subsection{Performance and memory footprint}

Across representative dataset sizes (spanning small, medium, and large scales), PyFgsea consistently reduced runtime while using substantially less memory than alternative Python implementations (Supplementary Table~S4). \rev{In the large benchmark (20k genes, 5k gene sets), PyFgsea completed in 1.35~s with 68.0~MB PeakRSS, compared with 18.00~s and 611.4~MB for BlitzGSEA and 189.42~s and 5891.1~MB for GSEApy.} These gains are primarily attributable to the Rust backend, pathway-level parallelism, and careful reuse of intermediate buffers.

\begin{figure}[H]
\centering
\includegraphics[width=0.88\textwidth]{assets/figure1_agreement.png}
\caption{Agreement of normalized enrichment scores (NES) across implementations on matched synthetic benchmark inputs. Panels compare PyFgsea with (A) GSEApy, (B) BlitzGSEA, and (C) R fgseaMultilevel. Dashed lines indicate identity; insets report concordance statistics. PyFgsea is near-identical to the R reference.}
\label{fig:agreement}
\end{figure}

\subsection{Statistical faithfulness to fgseaMultilevel}

To assess equivalence to the reference implementation, we compared enrichment scores (ES) and normalized enrichment scores (NES) against R fgseaMultilevel~\cite{ref4}. \rev{We observed near-identical ES/NES across a wide range of pathways (Figure~\ref{fig:agreement}), achieving rigorous quantitative concordance with the R reference (Pearson $r=1.000$; Spearman $\rho=1.000$; RMSE $=0.0135$; median absolute difference $|\Delta|=0.0070$ for NES; Figure~1C). ES values were numerically identical within machine precision. Transformed nominal $P$-values remained overall concordant with the reference implementation (overall Pearson $r=0.997$; Spearman $\rho=0.988$; RMSE $=0.1696$; median $|\Delta|=0.0467$ on $-\log_{10}P$), although sparse local trajectory windows could show visibly larger dispersion than ES or NES (worst-window Pearson $r=0.941$ on $-\log_{10}P$; Supplementary Table~S7). In such sparse local windows, ES and NES should therefore be regarded as the primary cross-implementation comparanda, whereas transformed nominal-$P$ agreement is more sensitive to tail-estimation noise and pathway-count filtering.} For multilevel $P$-values, PyFgsea employs a stricter default floor (\texttt{eps=1e-50}) compared to R's default (\texttt{eps=1e-10}), allowing for higher resolution in extreme tails. \rev{When \texttt{eps} is matched, deep-tail $P$-value differences are generally modest, but they become more apparent when few pathways pass size filters or when strongly enriched local windows occupy sensitive rare-event tails (Supplementary Figures~S3 and S4). PyFgsea thus preserves ES within machine precision and NES near-identically relative to the reference procedure while offering extended precision for rare events.}

\subsection{Trajectory case study: rolling-window enrichment}

PyFgsea enables rolling-window enrichment along pseudotime for large single-cell datasets. \rev{In a real GSE155254-derived HSC-like-to-erythroid trajectory, windowed enrichment identified temporally localized pathway activation patterns while maintaining feasible runtimes over hundreds to thousands of windows (Figure~\ref{fig:trajectory}). Across rolling-window benchmarks, the stateful runner yielded approximately 1.9-fold speedup in a conservative end-to-end stress test and up to 7.47-fold acceleration in a narrower 100-window component benchmark. Figure~\ref{fig:trajectory} uses the default full-window workflow without the optional terminal-truncation extension, which is evaluated separately in Supplementary Figure~S14.}

\begin{figure}[H]
\centering
\includegraphics[width=0.95\textwidth]{assets/figure2_real_trajectory.png}
\caption{Rolling-window GSEA along a GSE155254-derived HSC-like-to-erythroid trajectory. (A) UMAP colored by diffusion pseudotime. (B) Marker-defined states. (C) Smoothed HBB and MKI67 expression. (D) Window-wise NES for representative pathways; filled points mark within-window FDR $< 0.05$ after Benjamini--Hochberg correction across pathways. The default full-window workflow is shown; terminal-truncation is audited in Supplementary Figure~S14.}
\label{fig:trajectory}
\end{figure}
